\documentclass[11pt,a4paper]{article}
% Wspólny szablon dla kart supermocy (jedna karta = jedna strona A4)
% Użycie w karcie: \documentclass[11pt,a4paper]{article} % Wspólny szablon dla kart supermocy (jedna karta = jedna strona A4)
% Użycie w karcie: \documentclass[11pt,a4paper]{article} % Wspólny szablon dla kart supermocy (jedna karta = jedna strona A4)
% Użycie w karcie: \documentclass[11pt,a4paper]{article} \input{../preambula} \begin{document}

% Silnik tectonic (XeTeX): polskie znaki przez fontspec, nie przez inputenc/T1
\usepackage{fontspec}
\setmainfont{Arial}
\usepackage[polish]{babel}
\usepackage[a4paper,margin=1.5cm]{geometry}
\usepackage{amsmath,amssymb}
\usepackage{tikz}
\usetikzlibrary{arrows.meta,positioning,calc,shapes.geometric,decorations.pathreplacing}
\usepackage[most]{tcolorbox}
\usepackage{enumitem}
\usepackage{booktabs}
\usepackage{xcolor}

\pagestyle{empty}
\setlength{\parindent}{0pt}
\setlength{\parskip}{0.4em}

\definecolor{cel}{HTML}{1F4E79}
\definecolor{regula}{HTML}{E8F1FA}
\definecolor{pulapka}{HTML}{FDECEA}
\definecolor{okres}{HTML}{C0392B}
\definecolor{skonczony}{HTML}{1E8449}

% Ramka na regułę
\newtcolorbox{regula}[1][]{
  colback=regula, colframe=cel, boxrule=0.8pt, arc=2pt,
  left=6pt, right=6pt, top=4pt, bottom=4pt, title=#1,
  fonttitle=\bfseries, coltitle=white
}

% Ramka na pułapkę z kartkówki
\newtcolorbox{pulapka}[1][]{
  colback=pulapka, colframe=okres, boxrule=0.8pt, arc=2pt,
  left=6pt, right=6pt, top=4pt, bottom=4pt, title=#1,
  fonttitle=\bfseries, coltitle=white
}

% Nagłówek karty: numer, tytuł, cel
\newcommand{\naglowekKarty}[3]{%
  {\LARGE\bfseries\color{cel} Karta #1: #2}\\[2pt]
  {\large Cel: #3}\par
  \vspace{0.3em}\noindent\rule{\linewidth}{1pt}\par
}

% Pole na obliczenia ucznia: ramka z delikatną kratką, argument = wysokość
\newcommand{\miejsce}[1]{%
  \par\vspace{0.2em}\noindent
  \begin{tikzpicture}
    \draw[step=0.5cm, gray!25, very thin] (0,0) grid (\textwidth,#1);
    \draw[gray!60, rounded corners=2pt] (0,0) rectangle (\textwidth,#1);
    \node[anchor=north west, font=\scriptsize, text=gray] at (0.05,#1) {obliczenia};
  \end{tikzpicture}\par}
 \begin{document}

% Silnik tectonic (XeTeX): polskie znaki przez fontspec, nie przez inputenc/T1
\usepackage{fontspec}
\setmainfont{Arial}
\usepackage[polish]{babel}
\usepackage[a4paper,margin=1.5cm]{geometry}
\usepackage{amsmath,amssymb}
\usepackage{tikz}
\usetikzlibrary{arrows.meta,positioning,calc,shapes.geometric,decorations.pathreplacing}
\usepackage[most]{tcolorbox}
\usepackage{enumitem}
\usepackage{booktabs}
\usepackage{xcolor}

\pagestyle{empty}
\setlength{\parindent}{0pt}
\setlength{\parskip}{0.4em}

\definecolor{cel}{HTML}{1F4E79}
\definecolor{regula}{HTML}{E8F1FA}
\definecolor{pulapka}{HTML}{FDECEA}
\definecolor{okres}{HTML}{C0392B}
\definecolor{skonczony}{HTML}{1E8449}

% Ramka na regułę
\newtcolorbox{regula}[1][]{
  colback=regula, colframe=cel, boxrule=0.8pt, arc=2pt,
  left=6pt, right=6pt, top=4pt, bottom=4pt, title=#1,
  fonttitle=\bfseries, coltitle=white
}

% Ramka na pułapkę z kartkówki
\newtcolorbox{pulapka}[1][]{
  colback=pulapka, colframe=okres, boxrule=0.8pt, arc=2pt,
  left=6pt, right=6pt, top=4pt, bottom=4pt, title=#1,
  fonttitle=\bfseries, coltitle=white
}

% Nagłówek karty: numer, tytuł, cel
\newcommand{\naglowekKarty}[3]{%
  {\LARGE\bfseries\color{cel} Karta #1: #2}\\[2pt]
  {\large Cel: #3}\par
  \vspace{0.3em}\noindent\rule{\linewidth}{1pt}\par
}

% Pole na obliczenia ucznia: ramka z delikatną kratką, argument = wysokość
\newcommand{\miejsce}[1]{%
  \par\vspace{0.2em}\noindent
  \begin{tikzpicture}
    \draw[step=0.5cm, gray!25, very thin] (0,0) grid (\textwidth,#1);
    \draw[gray!60, rounded corners=2pt] (0,0) rectangle (\textwidth,#1);
    \node[anchor=north west, font=\scriptsize, text=gray] at (0.05,#1) {obliczenia};
  \end{tikzpicture}\par}
 \begin{document}

% Silnik tectonic (XeTeX): polskie znaki przez fontspec, nie przez inputenc/T1
\usepackage{fontspec}
\setmainfont{Arial}
\usepackage[polish]{babel}
\usepackage[a4paper,margin=1.5cm]{geometry}
\usepackage{amsmath,amssymb}
\usepackage{tikz}
\usetikzlibrary{arrows.meta,positioning,calc,shapes.geometric,decorations.pathreplacing}
\usepackage[most]{tcolorbox}
\usepackage{enumitem}
\usepackage{booktabs}
\usepackage{xcolor}

\pagestyle{empty}
\setlength{\parindent}{0pt}
\setlength{\parskip}{0.4em}

\definecolor{cel}{HTML}{1F4E79}
\definecolor{regula}{HTML}{E8F1FA}
\definecolor{pulapka}{HTML}{FDECEA}
\definecolor{okres}{HTML}{C0392B}
\definecolor{skonczony}{HTML}{1E8449}

% Ramka na regułę
\newtcolorbox{regula}[1][]{
  colback=regula, colframe=cel, boxrule=0.8pt, arc=2pt,
  left=6pt, right=6pt, top=4pt, bottom=4pt, title=#1,
  fonttitle=\bfseries, coltitle=white
}

% Ramka na pułapkę z kartkówki
\newtcolorbox{pulapka}[1][]{
  colback=pulapka, colframe=okres, boxrule=0.8pt, arc=2pt,
  left=6pt, right=6pt, top=4pt, bottom=4pt, title=#1,
  fonttitle=\bfseries, coltitle=white
}

% Nagłówek karty: numer, tytuł, cel
\newcommand{\naglowekKarty}[3]{%
  {\LARGE\bfseries\color{cel} Karta #1: #2}\\[2pt]
  {\large Cel: #3}\par
  \vspace{0.3em}\noindent\rule{\linewidth}{1pt}\par
}

% Pole na obliczenia ucznia: ramka z delikatną kratką, argument = wysokość
\newcommand{\miejsce}[1]{%
  \par\vspace{0.2em}\noindent
  \begin{tikzpicture}
    \draw[step=0.5cm, gray!25, very thin] (0,0) grid (\textwidth,#1);
    \draw[gray!60, rounded corners=2pt] (0,0) rectangle (\textwidth,#1);
    \node[anchor=north west, font=\scriptsize, text=gray] at (0.05,#1) {obliczenia};
  \end{tikzpicture}\par}


\begin{document}

\naglowekKarty{1}{Procent z liczby: klocki i ułamki}{obliczyć w pamięci dowolny procent z liczby}

\vspace{0.3em}
\begin{regula}[REGUŁA – ZBUDUJ PROCENT Z KLOCKÓW]
Procent to ułamek o mianowniku $100$: $\ p\%=\dfrac{p}{100}$. Nie mnóż w słupku, tylko składaj klocki:
\begin{itemize}[leftmargin=*, itemsep=2pt]
  \item $10\%$ to $\dfrac{1}{10}$: podziel przez $10$ (przesuń przecinek o $1$ w lewo). \quad $1\%$: podziel przez $100$.
  \item $5\%$ to \emph{połowa} z $10\%$. \quad $20\%$ to \emph{dwa razy} $10\%$.
  \item Przeliczniki: $50\%=\dfrac{1}{2}$, \ $25\%=\dfrac{1}{4}$, \ $75\%=\dfrac{3}{4}$, \ $20\%=\dfrac{1}{5}$, \ $5\%=\dfrac{1}{20}$.
\end{itemize}
\end{regula}

\vspace{0.4em}
\textbf{Krok po kroku: $35\%$ z $200$.} Najpierw klocek $10\%$, potem składamy:
\begin{enumerate}[leftmargin=*, itemsep=3pt]
  \item $10\%$ z $200=20$.
  \item $35\%=10\%+10\%+10\%+5\%$, czyli $3\cdot 20+10=70$ (bo $5\%$ to połowa z $20$).
\end{enumerate}
\textbf{Druga droga:} mnożenie. $35\%\cdot 200=\dfrac{35}{100}\cdot 200=35\cdot 2=70$. Oba sposoby muszą dać ten sam wynik.

\begin{center}
\begin{tikzpicture}[>=Stealth]
  % Pasek 200 podzielony na 10 klocków po 20 (każdy klocek = 10%)
  \foreach \i in {0,...,9} {
    \draw[thick] (\i*1.6,0) rectangle ++(1.6,0.8);
    \node[font=\small] at (\i*1.6+0.8,1.15) {$10\%$};
    \node[font=\small] at (\i*1.6+0.8,-0.35) {$20$};
  }
  \foreach \i in {0,1,2} { \fill[cel!35] (\i*1.6,0) rectangle ++(1.6,0.8); \draw[thick] (\i*1.6,0) rectangle ++(1.6,0.8); }
  \fill[okres!35] (4.8,0) rectangle ++(0.8,0.8);
  \draw[thick] (4.8,0) rectangle ++(1.6,0.8);
  \draw[dashed] (5.6,0) -- (5.6,0.8);
  \node[font=\small\bfseries, color=cel] at (2.4,0.4) {$3\cdot 20=60$};
  \node[font=\small\bfseries, color=okres] at (5.2,0.4) {$10$};
  \draw[decorate, decoration={brace, mirror, amplitude=5pt}, cel, thick] (0,-0.8) -- (5.6,-0.8)
    node[midway, below=6pt, font=\small\bfseries] {$35\%$ z $200=60+10=70$};
\end{tikzpicture}
\end{center}

\vspace{-0.3em}
\textbf{Przykłady rozstrzygnięte:}
\begin{center}
\renewcommand{\arraystretch}{1.5}
\begin{tabular}{@{}lll@{}}
\toprule
Zadanie & Klocki & Wynik \\
\midrule
$50\%$ z $80$ & połowa & $40$ \\
$75\%$ z $40$ & $\dfrac{1}{2}+\dfrac{1}{4}$: $20+10$ & $30$ \\[4pt]
$10\%$ z $230$ & przecinek w lewo & $23$ \\
$1\%$ z $450$ & przecinek o $2$ w lewo & $4{,}50$ \\
$15\%$ z $80$ & $10\%=8$, $5\%=4$ & $12$ \\
$7\%$ z $300$ & $1\%=3$, $7\cdot 3$ & $21$ \\
$120\%$ z $50$ & $100\%+20\%$: $50+10$ & $60$ \\
\bottomrule
\end{tabular}
\end{center}

\begin{pulapka}[PUŁAPKI Z KARTKÓWKI]
\begin{enumerate}[leftmargin=*, itemsep=2pt]
  \item $5\%\neq\dfrac{1}{5}$. Jedna piąta to $20\%$, a $5\%$ to $\dfrac{1}{20}$.
  \item $10\%$ z liczby to \emph{nie} „liczba minus $10$”. Dziel przez $10$.
  \item Więcej niż $100\%$ jest możliwe: $120\%$ z $50$ to więcej niż $50$.
\end{enumerate}
\end{pulapka}

\end{document}
